\cleardoublepage
\newrefsection
\chapter{外文翻译}

\section{预测性概念解码器}

语言模型有时会表现出与其自我报告相矛盾的行为。例如，一个被”越狱”的模型可能在提供制造炸弹的指令的同时，却声称自己并未意识到这一点。为了在不依赖模型”自述”的情况下理解其行为，我们提出了一种新的可解释性助手——预测性概念解码器（Predictive Concept Decoders, PCDs）。

\subsection{方法}

我们将可解释性重新框架为一个预测问题：如果一个系统真正理解模型的内部激活（activations），它应该能够预测该模型未来的行为。PCD 包含两个核心组件：
\begin{itemize}
    \item 编码器（Encoder）：将复杂庞大的模型激活状态压缩为稀疏的、人类可读的概念列表。
    \item 解码器（Decoder）：利用这些概念来回答关于模型行为的问题（例如”模型在想什么？“）。由于编码器看不见具体的问题，它被迫提取通用的概念，这使得其输出具有高度的可审计性。
\end{itemize}

\subsection{实验发现}

我们在模型持有未言明信息的场景中测试了 PCD，结果显示它能揭示模型试图隐藏的内容：
\begin{itemize}
    \item 越狱监测：当模型在被越狱状态下思考有害活动（如印假钞）但口头否认时，PCD 能准确解码出”如何印假钞”这一内部概念。
    \item 秘密提示与植入概念：当模型利用秘密提示作答，或被人为植入特定概念时，PCD 能将这些隐蔽的驱动因素语言化，而模型本身往往会编造虚假的理由来解释其行为。
\end{itemize}

\subsection{可扩展性与结论}

实验表明，PCD 具有良好的扩展性（Scaling laws）：随着训练数据的增加，编码器提取的概念变得更具可解释性，下游任务（如检测欺骗）的准确率也随之提升。这证明了通过端到端训练来构建可解释性工具的可行性，未来有望实现对大模型推理过程的完全透明化监控。

\section{小数迷思概念分类}

小数是数学学习中的基础难点，学生常持有持久的迷思概念（Misconceptions）。尽管已有大量分散的研究，但缺乏一个统一的分类体系来支持智能辅导系统（ITS）的开发。本文旨在通过整合过往文献，建立一个基于错误例题（Erroneous Examples）学习的小数迷思概念分类学。

\subsection{小数迷思概念分类体系 (Taxonomy)}

研究者构建了一个层级分类图，重点关注先验知识（如整数、分数知识）对小数理解的干扰：
\begin{itemize}
    \item Megz (Longer is Larger)：认为小数位数越多，数值越大。这是受整数规则（如 25 > 7）的负迁移影响，导致学生认为 0.25 > 0.7。
    \item Segz (Shorter is Larger)：认为小数位数越少，数值越大。这是受分数知识（分母越大数值越小）的负迁移影响，导致学生认为 0.2 > 0.25。
    \item Negz：将小于 1 的小数错误地视为负数。
    \item Regz：完全将小数视为自然数处理。
    \item Ovegz：对位值（place value）理解错误，例如认为 0.67 等于 67 个十分之一。
    \item Zegz：认为长小数更大，除非有前导零（如认为 .045 < .4 但 145 > .4）。
\end{itemize}

\subsection{应用与结论}

该分类学已被用于开发智能辅导系统（如 MathTutor），通过向学生展示包含特定错误的例题（Erroneous Examples），引导学生解释错误原因，从而修正其迷思概念。这种分类方法为个性化诊断和干预提供了结构化基础。

\section{儿童小数习得}

小数学习需要概念转变，即从自然数概念扩展到有理数。本研究旨在通过横断研究（3-6年级），区分小数长度（length）和数值（digit value）对儿童判断的独立影响，并探究儿童对”零”的作用的理解随年龄的变化。

\subsection{实验设计}

研究选取了 156 名 3 年级至 6 年级的学生，进行小数配对比较任务。题目设计包含多种类型：长度与数值一致（Congruent，如 0.20 vs 0.1）、长度与数值冲突（Incongruent，如 0.1 vs 0.02）以及涉及零的位置变化。

\subsection{核心发现}

\begin{itemize}
    \item 低年级（3-4年级）：自然数主导。 学生主要受数值影响，认为数字大的小数更大；同时受长度影响，倾向认为”长的小数更大”（符合 Megz 偏差）。他们常错误地认为在小数末尾加零会改变数值大小。
    \item 5年级：过渡与反向偏差。 学生正确率显著提升，但出现一种反向偏差：倾向认为”短的小数更大”。这可能是对自然数规则的过度矫正，或受分数知识干扰。他们对小数点后紧跟的零（placeholder zero）理解较好，但仍对末尾的零感到困惑。
    \item 6年级：概念掌握。 大多数学生不再受长度或数值的显著干扰，能够正确理解零在不同位置的作用，表现出专家水平。
\end{itemize}

\subsection{结论}

小数概念的发展是缓慢且渐进的，存在中间状态（如5年级的表现）。聚类分析显示，低年级学生多属于仅关注数值或受长度影响的群体，而高年级学生则逐渐过渡到能综合处理信息的专家组。

\section{自然数偏差}

人类在比较比率（分数、概率）时常表现出”自然数偏差”（偏好分子更大的选项，即便比率相同或更小）。学界对此有两种解释：一种是策略性假说（仅在任务困难时退回使用自然数比较），另一种是内在性假说（自然数比较是比率处理过程中不可避免的内在部分）。本研究旨在通过非符号化任务验证这两种假说。

\subsection{实验方法}

\begin{itemize}
    \item 实验 1：赌博任务。 向受试者展示两个装有不同颜色球的罐子图片（非符号化），要求选择获胜概率大的罐子。在概率相等（如 50\% vs 50\%）但球数不同（如 4/8 vs 1/2）的情况下，受试者显著偏好球数更多（分子更大）的选项，表现出明显的自然数偏差。
    \item 实验 2：快速概率判断与贝叶斯建模。 通过大量试次测试受试者在不同比率距离下的反应，并构建计算模型（如分母忽略模型、整体比率模型、策略性模型、内在偏差模型）来拟合数据。
\end{itemize}

\subsection{核心发现}

贝叶斯模型比较结果显示，内在自然数偏差模型（Intrinsic Whole Number Bias）最能解释人类行为。这意味着：
\begin{itemize}
    \item 偏差并非因为人们在困难时”放弃”了比率计算。
    \item 相反，人们在进行比率比较时，会同时计算并加权考虑”整体比率”和”分子的自然数数值”。
    \item 即使受试者能准确知觉比率，自然数的大小仍会作为一种内在的、自动激活的量级信息干扰决策。
\end{itemize}

\subsection{结论}

自然数偏差不是简单的策略失误，而是人类认知系统处理比例信息时的一种根深蒂固的特征，可能源于早期发展的非言语数值表征系统。

\section{自然数偏差的变异性：对象、时机、方式与原因}

本文是一篇理论评论，旨在解释为何自然数偏差（NNB）的表现存在巨大的变异性（Variability）。并非所有人在所有任务中都会表现出偏差，同一个人在不同试次中也可能表现不同。作者提出了动态策略选择（Dynamic Strategy Choice）框架来解释这一现象。

\subsection{变异性的来源}

\begin{itemize}
    \item 谁与何时：偏差不仅存在于儿童中，也存在于成人甚至数学专家中（在快速反应任务中）。任务类型（如分数比较、代数表达式判断）显著影响偏差的激活程度。
    \item 如何产生：作者反驳了单纯的”双重加工理论”（直觉 vs 分析），认为更应关注表征的强度与精确度。偏差的产生是因为自然数的大小表征（magnitude representations）或运算模式比有理数的表征更强、更精确、且更容易被激活。
\end{itemize}

\subsection{动态策略选择模型}

人们在解决问题时，会根据当前任务的启示性（affordances）和自身知识库，动态地选择策略。
\begin{itemize}
    \item 当有理数表征（如分数的数轴位置）不够强时，个体会退而求其次，依赖更稳固的自然数表征（如比较分子大小），从而导致偏差。
    \item 即便在使用费力的分析性策略时，如果自然数知识被过度激活，也会导致错误。
\end{itemize}

\subsection{结论}

克服自然数偏差的关键不在于单纯抑制直觉，而在于通过练习和经验，提高有理数表征的强度和精确度，使其在策略竞争中能够胜过自然数表征。
